%% kapa: exam=1 round=26 year=2015 period=2 form=A subject=appraisal_law subject_ko=부동산관계법규 seq=27 number=27
%% source: Q-Net 감정평가사 기출문제 (cst00309) artlSeq=5138001 file=extracted/1008/5138001/제26회_감정평가사 1차_기출/제26회_감정평가사 1차_기출/2교시_감평_A형.pdf
%% conversion: kordoc 4.24.51 -> tools/gen1.py (auto); review=proofread-sonnet; flags=none
\begin{question}{27}
국유재산법령상 국유재산의 관리에 관한 설명으로 옳지 않은 것은?
\begin{choices}
  \item 국유재산의 사용료가 납부기한까지 납부되지 아니한 경우 연체료 부과대상이 되는 연체기간은 납기일부터 60개월을 초과할 수 없다.
  \item 국가는 과오납된 국유재산의 사용료를 반환하는 경우에는 과오납된 날의 다음 날부터 반환하는 날까지의 기간에 대하여 이자를 가산하여 반환한다.
  \item 지방자치단체가 은닉된 국유재산이나 소유자 없는 부동산을 발견하여 신고한 경우에는 그 재산가격의 2분의 1의 범위에서 그 지방자치단체에 국유재산을 양여하거나 보상금을 지급할 수 있다.
  \item 정당한 사유 없이 국유재산을 점유하거나 이에 시설물을 설치한 경우「행정대집행법」을 준용하여 철거할 수 없다.
  \item 은닉된 국유재산을 선의(善意)로 취득한 후 그 재산을 국가에 자진 반환한 자에게 같은 재산을 매각하는 경우에는 그 매각대금을 이자 없이 12년 이하에 걸쳐 나누어 내게 할 수 있다.
\end{choices}
\end{question}
\answer{4}
